\documentclass[letterpaper]{article}
\usepackage{iftex}
\ifPDFTeX
  % 한글 미리보기 빌드에는 XeLaTeX 또는 LuaLaTeX를 사용한다.
  \errmessage{Korean version requires XeLaTeX or LuaLaTeX}
\else
  % AAAI 지면 및 2단 편집 형식을 따르는 휴대용 국문 미리보기.
  % 투고본이 아니라 국문 검토용이다. 정본은
  % 2026-07-26_embodied_reasoning_results.tex 이며, 수치는 전부 그 파일과 일치시킨다.
  \usepackage[margin=0.75in,columnsep=0.375in]{geometry}
  \usepackage{authblk}
  \setlength{\affilsep}{0pt}
  \twocolumn
  \usepackage{fontspec}
  % Noto Sans CJK KR에는 이탤릭·스몰캡스 자형이 없다. 기울임은 합성으로 채우고,
  % 스몰캡스가 필요한 자리는 아래 \scmark(대문자)로 대체한다.
  \IfFontExistsTF{Noto Sans CJK KR}{
    \setmainfont[AutoFakeBold=2,AutoFakeSlant=0.2]{Noto Sans CJK KR}
    \setsansfont[AutoFakeBold=2,AutoFakeSlant=0.2]{Noto Sans CJK KR}
  }{
    \IfFontExistsTF{NanumMyeongjo}{
      \setmainfont{NanumMyeongjo}
      \setsansfont{NanumGothic}
    }{
      \setmainfont{DejaVu Sans}
      \setsansfont{DejaVu Sans}
    }
  }
  \XeTeXlinebreaklocale "ko"
  \XeTeXlinebreakskip=0pt plus 1pt
\fi
\usepackage[hyphens]{url}
\usepackage{graphicx}
\urlstyle{rm}
\def\UrlFont{\rm}
\usepackage{natbib}
\usepackage{caption}
\renewcommand{\figurename}{그림}
\renewcommand{\tablename}{표}
\frenchspacing

\usepackage{booktabs}
\usepackage{amsmath}
\usepackage{amssymb}
\usepackage{xcolor}
\usepackage{tikz}

% 초안 색상 규약:
%   파랑 = 측정값, 빨강 = 최종 결과로 기각된 사전등록 주장.
\definecolor{measuredblue}{RGB}{0,73,160}
\definecolor{pendingred}{RGB}{190,35,35}
\newcommand{\obs}[1]{\textcolor{measuredblue}{#1}}
\newcommand{\rejected}[1]{\textcolor{pendingred}{#1}}
\newcommand{\pp}{\,\mathrm{pp}}
\newcommand{\method}{\textsc{EGO}}
\newcommand{\cmark}{\checkmark}
% 국문 폰트에 스몰캡스가 없으므로 규칙 트리거 어휘는 대문자로 표시한다.
\newcommand{\scmark}[1]{\MakeUppercase{#1}}

\setcounter{secnumdepth}{0}
\title{실험 초안: 근거 기반 선택으로서의 체화 추론}
\author{익명 투고}
\affil{익명 소속}
\date{2026년 7월 26일 \\ \small 국문 검토본 --- 정본은 영문
  \texttt{2026-07-26\_embodied\_reasoning\_results.tex}}

\begin{document}
\setlength{\emergencystretch}{1em}
\maketitle

\section{실험}
\subsection{체화 추론}
\label{sec:er-capability}

\paragraph{평가 프로토콜과 증거 제시 순서.}
Ego4D GoalStep 홀드아웃 분할의 예측 시점에서 평가한다.
동결된 월드 모델은 관측과 완료 행동 이력을 입력받은 뒤 순서를 무작위로 섞은
Top-$K$ 행동 경계를 제안하며, 월드 모델의 점수와 순위는 정책에 공개하지 않는다.
별도 언급이 없는 한, 모든 비교는 동일한 1인칭 평가 코호트와 공통 커버 집합
($a_t^{\mathrm{GT}}\!\in D_t$; \obs{$n{=}1{,}000$})을 사용한다.
증거는 결과에서 메커니즘으로 이어지는 순서로 구성한다. 먼저 통제된 선택 비교를
제시하고, 다음으로 역량을 분해하며, 마지막으로 이력과 과업 믿음에 대한 쌍별
개입을 수행한다. 이 순서는 \emph{무엇을 선택하는가}의 개선과
\emph{추론 흔적을 어떻게 작성하는가}의 피상적 변화를 구분한다.

\subsubsection{정답 전용 감독이 아니라 후보 정렬이 선택을 개선한다}
\label{sec:er-selection}

주요 비교에서는 백본, 예시, 학습 예산, 평가 프롬프트를 통제하며, 유일한 차이는
학습 중 월드 모델 후보 집합을 노출하는지 여부이다. 후보 정렬
($\theta_{\mathrm{CE}}$)의 선택 정확도는 \obs{$28.8\%$}로, 정답 전용 감독
(GT-only)의 \obs{$21.9\%$}보다 높다. 형식 오류를 제외한 공식 쌍별 집합
(\obs{$n{=}948$}, 비디오 클러스터 85개)에서 두 조건의 차이는
\obs{$+7.70\pp$}이며, 비디오 클러스터 부트스트랩
\obs{$95\%$ 신뢰구간은 $[+3.73,+11.93]$}이다.
반면 GT-only와 비정렬 정책의 차이는 \obs{$+0.9\pp$}에 불과하다
(\obs{$95\%$ 신뢰구간 $[-2.0,+3.7]$}). 후보 정렬은 동결된 월드 모델의
Top-\obs{$1$} 정책도 넘어선다(\obs{$28.8\%$} 대 \obs{$24.2\%$}; 엄격 기준
\obs{$+4.6\pp$}). 공식 쌍별 게이트로는 \obs{$+6.06\pp$}이다
(\obs{$95\%$ 신뢰구간 $[+0.87,+11.22]$}; \obs{$n{=}957$}, 클러스터 86개).
따라서 이러한 향상은 정답을 추가로 노출했거나 월드 모델의
순위를 복사했기 때문으로 설명할 수 없다. 시각적으로 그럴듯한 대안들 사이를
판별하는 법을 학습해야만 얻을 수 있는 향상이다.
그림~\ref{fig:capability}b는 분기하는 학습 경로를 따라 SelAcc를 추적한다.
후보 정렬이 정책을 월드 모델 기준선 너머로 옮기는 단계이다.

\begin{figure*}[t]
\centering
\resizebox{0.96\textwidth}{!}{%
\begin{tikzpicture}[x=1cm,y=1cm,font=\small]
  % 패널 (a): 제목, 범례, 도표 영역을 분리한다.
  \node[anchor=west,font=\bfseries] at (0,6.15)
    {(a) 올바른 사전예측의 유지와 오류의 교정};
  \fill[measuredblue] (0.25,5.45) circle (2.5pt);
  \node[anchor=west] at (0.45,5.45) {$G_1$: WM 정답 유지};
  \fill[measuredblue!35] (4.15,5.45) circle (2.5pt);
  \node[anchor=west] at (4.35,5.45) {$G_2$: WM 오답 교정};

  \foreach \x/\lab in {1.5/0,2.75/10,4.0/20,5.25/30,6.5/40}{
    \draw[gray!30] (\x,0.85) -- (\x,4.65);
    \node[anchor=north,text=measuredblue] at (\x,0.72) {\lab};
  }
  \node[anchor=north] at (4.0,0.15) {조건부 선택 정확도 (\%)};
  \node[anchor=east] at (1.15,4.45) {기본};
  \node[anchor=east] at (1.15,3.40) {GT-only};
  \node[anchor=east] at (1.15,2.35) {후보-CE};
  \node[anchor=east] at (1.15,1.30) {\method{}};

  \foreach \y/\gOne/\gTwo in {
    4.45/5.2125/3.8375,
    3.40/4.8125/4.10,
    2.35/6.175/4.975,
    1.30/6.8125/4.70}{
    \draw[line width=2.2pt,gray!45] (\gTwo,\y)--(\gOne,\y);
    \fill[measuredblue!35] (\gTwo,\y) circle (2.5pt);
    \fill[measuredblue] (\gOne,\y) circle (2.5pt);
  }
  \node[above=2pt,text=measuredblue] at (5.2125,4.45) {\obs{$29.7$}};
  \node[below=2pt,text=measuredblue] at (3.8375,4.45) {\obs{$18.7$}};
  \node[above=2pt,text=measuredblue] at (4.8125,3.40) {\obs{$26.5$}};
  \node[below=2pt,text=measuredblue] at (4.10,3.40) {\obs{$20.8$}};
  \node[above=2pt,text=measuredblue] at (6.175,2.35) {\obs{$37.4$}};
  \node[below=2pt,text=measuredblue] at (4.975,2.35) {\obs{$27.8$}};
  \node[above=2pt,text=measuredblue] at (6.8125,1.30) {\obs{$42.5$}};
  \node[below=2pt,text=measuredblue] at (4.70,1.30) {\obs{$25.6$}};

  % 패널 (b): 학습 스텝 대비 SelAcc. 색 계열은 패널 (a)와 공유한다.
  % x는 옵티마이저 스텝 0--900을 1.05--7.15cm로 사상한다: x(s) = 1.05 + s*0.0067778
  %   (s=522 -> 4.588 = theta_CE 인계 지점, s=890 -> 7.082 = Retrospection 종료).
  % y는 19--30%를 0.85--4.65cm로 사상한다: y(v) = 0.85 + (v-19)*0.345455.
  % 마커는 실제 평가된 상태만 표시하며, 곡선은 그 사이를 보간한다.
  \begin{scope}[xshift=9.05cm]
    \node[anchor=west,font=\bfseries] at (0,6.15)
      {(b) 학습 경로에 따른 선택 정확도};

    \begin{scope}[font=\footnotesize]
      \draw[black!45,dotted,line width=1.2pt] (0.88,5.45) -- (1.26,5.45);
      \node[anchor=west] at (1.32,5.45) {기본};
      \draw[measuredblue!45,line width=1.5pt] (2.16,5.45) -- (2.54,5.45);
      \fill[measuredblue!45] (2.35,5.45) circle (2.2pt);
      \node[anchor=west] at (2.60,5.45) {GT-only};
      \draw[measuredblue,line width=1.7pt] (3.94,5.45) -- (4.32,5.45);
      \fill[measuredblue] (4.13,5.45) circle (2.4pt);
      \node[anchor=west] at (4.38,5.45) {후보-CE};
      \draw[measuredblue!70!black,line width=2.3pt] (5.90,5.45) -- (6.28,5.45);
      \fill[measuredblue!70!black] (6.09,5.45) circle (2.6pt);
      \node[anchor=west] at (6.34,5.45) {\method{}};
    \end{scope}

    \foreach \y/\lab in {1.1955/20,2.9227/25,4.65/30}{
      \draw[gray!30] (1.05,\y) -- (7.15,\y);
      \node[anchor=east,text=measuredblue] at (0.92,\y) {\lab};
    }
    \node[rotate=90,anchor=south] at (0.08,2.75)
      {엄격 SelAcc@Top-$K$ (\%, $\uparrow$)};
    \foreach \x/\lab in {1.05/0,2.744/250,4.439/500,6.133/750}{
      \node[anchor=north,text=measuredblue] at (\x,0.72) {\lab};
    }
    \node[anchor=north] at (4.10,0.08)
      {Prospection $\rightarrow$ Retrospection 학습 스텝};

    % 두 Retrospection 실행은 동일한 theta_CE 어댑터에서 분기한다.
    \draw[gray!45,dashed,line width=0.7pt] (4.588,0.85) -- (4.588,3.85);
    \node[anchor=south west,font=\scriptsize,text=black!55,fill=white,
          fill opacity=0.85,text opacity=1,inner sep=1pt]
      at (4.66,0.92) {$\theta_{\mathrm{CE}}$ 인계};

    % 동결 월드 모델 Top-1 기준선.
    \draw[black!60,dashed,line width=1.0pt] (1.05,2.6464) -- (7.15,2.6464);
    \node[anchor=north east,font=\footnotesize,fill=white,fill opacity=0.86,
          text opacity=1,inner sep=1.5pt] at (7.12,2.60)
      {동결 WM Top-1 $L_0=\obs{24.2\%}$};

    % 기본: 백본 동결. 정의상 상수이며 모든 조건의 공통 출발점이다.
    \draw[black!45,dotted,line width=1.2pt] (1.05,1.5409) -- (7.15,1.5409);
    \fill[black!55] (1.05,1.5409) circle (2.6pt);
    \node[anchor=north west,inner sep=2pt,text=black!65,fill=white,
          fill opacity=0.85,text opacity=1] at (1.00,1.49) {기본 \obs{$21.0$}};

    % GT-only: Prospection 구간의 정답 전용 선택 CE.
    \draw[measuredblue!45,line width=1.7pt]
      (1.05,1.5409) .. controls (2.10,1.70) and (3.40,1.82) .. (4.588,1.8518);
    \fill[measuredblue!45] (4.588,1.8518) circle (2.8pt);
    \node[anchor=south east,inner sep=2pt,text=measuredblue]
      at (4.54,1.90) {GT-only \obs{$21.9$}};

    % 후보-CE: 동일 구간의 후보 정렬 선택 CE.
    \draw[measuredblue,line width=1.9pt]
      (1.05,1.5409) .. controls (2.10,3.30) and (3.40,4.05) .. (4.588,4.2455);
    \fill[measuredblue] (4.588,4.2455) circle (2.9pt);
    \node[above=3pt,text=measuredblue] at (4.30,4.2455) {후보-CE \obs{$28.8$}};
    \node[font=\scriptsize,fill=white,inner sep=1pt,text=measuredblue]
      at (1.92,3.62) {\obs{$+7.8\pp$}};

    % EGO: 재생 및 근거화 오버샘플을 적용한 Retrospection.
    \draw[measuredblue!70!black,line width=2.5pt]
      (4.588,4.2455) -- (7.082,4.2455);
    \fill[measuredblue!70!black] (7.082,4.2455) circle (3.1pt);
    \node[anchor=south east,inner sep=2pt,text=measuredblue]
      at (7.14,4.31) {\method{} \obs{$28.8$}};
  \end{scope}
\end{tikzpicture}
}
\caption{\textbf{후보 정렬은 정답 사전분포가 아니라 근거 기반 선택을 학습시킨다.}
(a) $G_1$은 월드 모델 Top-\obs{$1$}이 정답일 때의 유지를 평가하고,
$G_2$는 오답일 때의 교정을 평가한다. 후보 정렬은 둘 모두를 개선하지만,
GT-only는 정답을 보았음에도 $G_1$을 저하시킨다. 최종 \method{} 모델의 유지율은
\obs{$42.5\%$}에 이른다. (b) 분기하는 학습 경로에서 학습 스텝 대비 엄격
SelAcc@Top-$K$를 표시한다. 모든 조건은 동일한 동결 출발점(\obs{$21.0\%$})을
공유한다. Prospection 구간(선택 CE \obs{$522$} 스텝)에서 후보 정렬은
\obs{$28.8\%$}에 도달하는 반면 정답 전용 대조군은 \obs{$21.9\%$}에 그친다.
이어 \method{}는 $\theta_{\mathrm{CE}}$ 어댑터에서 분기해
Retrospection-with-Replay($\rho{=}0.15$)를 \obs{$368$} 스텝 더 학습하며
\obs{$28.8\%$}를 유지한다. 기본 모델은 동결되어 있으므로 정의상 상수이다.
선은 시각적 안내선이다. 후보 정렬이 정책을 점선으로 표시한 동결 WM Top-1 기준
$L_0=\obs{24.2\%}$ 너머로 옮기는 단계이며(\obs{$+7.8\pp$}),
Retrospection은 그 위치를 보존한다.}
\label{fig:capability}
\end{figure*}

\subsubsection{향상이 나타나는 지점: 다섯 가지 역량 축}
\label{sec:er-axes}

집계 정확도만으로는 정책이 올바른 예측 사전분포를 보존하는지, 아니면 잘못된
사전분포를 교정하는지 알 수 없다. 따라서 선택을 $G_1$ 유지(월드 모델
Top-\obs{$1$}이 정답)와 $G_2$ 교정(월드 모델 Top-\obs{$1$}이 오답)으로
나눈다. 그림~\ref{fig:capability}a와 표~\ref{tab:capability}에서 보듯,
후보 정렬은 $G_1$을 \obs{$29.7\%$}에서 \obs{$37.4\%$}로, $G_2$를
\obs{$18.7\%$}에서 \obs{$27.8\%$}로 높인다. 반면 GT-only는 $G_2$를
\obs{$20.8\%$}로 소폭 높이지만 $G_1$은 \obs{$26.5\%$}로 낮춘다.
최종 \method{} 모델은 $G_1$을 \obs{$42.5\%$}까지 더 높이는 한편
$G_2$는 \obs{$25.6\%$}에 자리 잡는다.
따라서 추론 시점에 예측 경계를 제공하는 것만으로는 충분하지 않다. 정책은 그
경계의 선두 가설을 언제 유지하고 언제 기각해야 하는지를 학습해야 한다.

\begin{table*}[t]
\centering
\small
\setlength{\tabcolsep}{5pt}
\begin{tabular}{p{0.245\textwidth}ccccc}
\toprule
역량 지표 & 기본 & GT-only & 후보-CE & \method{} & 평가 대상 \\
\midrule
SelAcc@Top-$K$ (\%, $\uparrow$)
  & \obs{$21.0$} & \obs{$21.9$} & \textbf{\obs{$28.8$}}
  & \textbf{\obs{$28.8$}} & 전체 근거 기반 선택 \\
$G_1$: WM 정답 유지 (\%, $\uparrow$)
  & \obs{$29.7$} & \obs{$26.5$} & \obs{$37.4$}
  & \textbf{\obs{$42.5$}} & 유효한 예측 사전분포 보존 \\
$G_2$: WM 오답 교정 (\%, $\uparrow$)
  & \obs{$18.7$} & \obs{$20.8$} & \textbf{\obs{$27.8$}}
  & \obs{$25.6$} & WM Top-\obs{$1$} 이외 재선택 \\
연속 행동 재현율 (\%, $\uparrow$)
  & \obs{$28.4$} & \obs{$32.9$} & \obs{$43.5$}
  & \textbf{\obs{$48.4$}} & 궤적과 일관된 행동 복원 \\
연속 행동 정밀도 (\%)
  & \obs{$72.7$} & \obs{$71.3$} & \obs{$71.1$}
  & \obs{$68.2$} & 신호 신뢰도(통제) \\
증거 유용성 (정확도 \(\Delta\), pp, $\uparrow$)
  & \obs{$+13.5$} & \obs{$+3.2$} & \obs{$+12.3$}
  & \textbf{\obs{$+18.3$}} & 증거--의사결정 결합 \\
증거 언급률 (\%)
  & \obs{$22.4$} & \obs{$37.8$} & \obs{$32.4$}
  & \obs{$17.5$} & 표면 빈도(통제) \\
믿음--행동 반복 (\%, $\downarrow$)
  & \obs{$0.0$} & \obs{$28.8$} & \obs{$2.1$}
  & \textbf{\obs{$0.0$}} & 행동 복사를 넘어선 믿음 내용 \\
\bottomrule
\end{tabular}
\caption{\textbf{공통 커버 집합의 역량 프로파일}
(\obs{$n{=}1{,}000$}). 각 열은 차례로 비정렬 정책, 정답 전용 대조군,
후보 정렬 정책, 최종 \method{} 모델에 대응한다.
연속 행동은 집합의 \obs{$31.0\%$}
(\obs{$310/1{,}000$})를 이룬다. 안정적인 정밀도는 주된 한계가 더 잡음이 많은
연속 행동 추측이 아니라 복원임을 보여준다. 증거 유용성은 이전 행동 패턴을
명시적으로 활용한 추론 흔적과 그렇지 않은 추론 흔적 사이의 SelAcc 차이로,
인과 효과가 아닌 관찰적 지표이다. 월드 모델 후보 풀은 홀드아웃 예측 시점의
\obs{$43.4\%$}를 포괄하므로, \obs{$28.8\%$}의 전체 집합 환산값은
\obs{$12.5\%$}이며 WM Top-\obs{$1$} 정책의 \obs{$10.5\%$}와 대비된다.
최종 \method{} 체크포인트도 같은 코호트에서
측정한다. SelAcc에는 엄격한 분모를 사용하고, 조건부 역량 축에서는 위 정의에
따라 형식이 잘못된 추론 흔적을 제외한다.}
\label{tab:capability}
\end{table*}

궤적 연속성에서도 동일한 패턴이 나타난다. 바로 앞의 행동을 이어가는 정답 행동은
평가 집합의 \obs{$31.0\%$}를 차지한다. 후보 정렬은 그중 \obs{$43.5\%$}를
복원하는 반면, 기본 모델은 \obs{$28.4\%$}, GT-only는 \obs{$32.9\%$}를
복원하며, 최종 \method{} 모델은 \obs{$48.4\%$}에 이른다.
한편 연속 행동 정밀도는 평가한 모든 모델 상태에서 \obs{$68$--$73\%$} 범위에
머문다. 정밀도가 비교적 안정적이라는 사실은 이전 행동을 더 느슨하게 반복하는
경향으로 결과를 설명할 수 없음을 뜻한다. 정렬은 주로 고정밀 시간 단서의
재현율을 개선한다.

\subsubsection{인과 검정: 학습된 이점은 행동 이력을 통해 전달된다}
\label{sec:er-history}

다음으로 이미지, 후보 집합, 체크포인트는 고정한 채 완료 행동 이력을 제거한다.
후보 정렬의 성능은 \obs{$10.1\pp$} 감소한다
(\obs{$28.8\%\!\rightarrow18.7\%$};
\obs{$95\%$ 신뢰구간 $[+6.8,+13.3]$}). 반면 기본 모델과 GT-only의 감소폭은
각각 \obs{$2.9\pp$}와 \obs{$4.5\pp$}이고, 재생 체크포인트는 \obs{$6.4\pp$}
감소한다. 이 개입은 최종 체크포인트에서는 재실행하지 않았다.
더 중요한 점은 이력이 없을 때
후보 정렬이 기본 모델(\obs{$+0.6\pp$}, \obs{$95\%$ 신뢰구간
$[-1.8,+3.0]$}) 및 GT-only(\obs{$+1.3\pp$}, \obs{$95\%$ 신뢰구간
$[-1.2,+3.9]$})와 구분되지 않는다는 것이다. 이 쌍별 추정값은 이러한 붕괴를
정량화한다. 학습이 단순히 더 나은 행동 사전분포를 주입했다면
이력을 가린 뒤에도 이점이 남아야 한다. 그러나 이점은 사라지며, 이는 후보 정렬이
모델에 궤적을 읽고 이를 활용하여 월드 모델 제안 중 하나를 선택하도록
학습시킨다는 것을 보여준다.

\subsubsection{더 많은 증거가 아니라 유용한 증거}
\label{sec:er-evidence}

마지막으로 생성된 추론 흔적이 의사결정과 관련된 상태를 포함하는지, 아니면
설득력 있는 표현만 담는지를 묻는다. GT-only 감독에서는 이전 행동 패턴을
언급하는 빈도가 높아지지만(\obs{$22.4\%\!\rightarrow37.8\%$}), 그러한 언급과
연관된 정확도 격차는 \obs{$+13.5\pp$}에서 \obs{$+3.2\pp$}로 붕괴한다.
후보 정렬도 언급 빈도를 \obs{$32.4\%$}로 높이지만, 훨씬 큰
\obs{$+12.3\pp$}의 증거 유용성을 유지한다. 최종 \method{} 체크포인트는
그러한 증거를 오히려 덜 언급하지만(\obs{$17.5\%$}), 조건부 유용성은
\obs{$+18.3\pp$}까지 높인다. 따라서 핵심적인 차이는 모델이
시간적 증거를 사용했다고 \emph{말하는지}가 아니라, 인용한 증거가 올바른
의사결정과 잘못된 의사결정을 구분하는지에 있다.

\begin{table}[t]
\centering
\scriptsize
\setlength{\tabcolsep}{2.5pt}
\begin{tabular}{lcccc}
\toprule
메커니즘 진단 & 기본 & GT-only & 후보-CE & \method{} \\
\midrule
증거 언급 (\%) & \obs{$22.4$} & \obs{$37.8$} & \obs{$32.4$} & \obs{$17.5$} \\
언급 시 정확도 (\%) & \obs{$31.8$} & \obs{$24.1$} & \obs{$38.4$} & \obs{$44.7$} \\
미언급 시 정확도 (\%) & \obs{$18.4$} & \obs{$21.0$} & \obs{$26.1$} & \obs{$26.4$} \\
증거 유용성 (pp) & \obs{$+13.5$} & \obs{$+3.2$} & \obs{$+12.3$} & \obs{$+18.3$} \\
\midrule
믿음--행동 반복 (\%) & \obs{$0.0$} & \obs{$28.8$} & \obs{$2.1$} & \obs{$0.0$} \\
믿음 교환 민감도 & \obs{$0.073$} & \obs{$0.085$} & \obs{$0.093$} & --- \\
\bottomrule
\end{tabular}
\caption{\textbf{표면적 및 개입적 진단.} 증거 유용성은 조건부 연관성이며
기술적으로 해석한다. 믿음 교환 민감도는 개입적 지표이다. 추론은 고정한 채
믿음 접두부만 교환하고, 재서술 통제 조건의 선택 변화율을 차감한다
(기본, GT-only, 후보-CE에 대해 \obs{$n{=}400$}).
최종 \method{} 모델에는 믿음 교환 평가를 수행하지 않았으므로, 직전
체크포인트의 값을 옮겨 적지 않고 빈칸으로 둔다.
후보 정렬에서 낮은 반복률은 믿음 필드가 선택된 행동을
단순히 복사한 것이 아님을 보여준다.}
\label{tab:mechanism}
\end{table}

믿음 채널은 보완적인 통제를 제공한다. 정답 전용 감독에서 과업 믿음은 추론 흔적의
\obs{$28.8\%$}에서 선택된 행동과 동일해진다. 기본 모델은 \obs{$0.0\%$},
후보 정렬 이후에는 \obs{$2.1\%$}이며, 최종 \method{} 모델도 반복률은
\obs{$0.0\%$}이다. 그러나 추론 텍스트를 고정한 채 믿음을
교환하면 후보 정렬 예시의 \obs{$9.5\%$}에서 선택된 행동이 바뀌며, 재서술
자기통제에서는 \obs{$0.25\%$}가 바뀐다. 이 통제값을 차감한 믿음 교환 민감도는
\obs{$9.3\%$}이다. 낮은 반복률과 0이 아닌
개입 민감도를 함께 고려하면, 믿음 필드는 행동의 재진술로 퇴화하지 않으면서
의사결정 관련 상태를 전달할 수 있다. 관찰적 증거 분할은 인과적으로 해석하지
않는다. 믿음 주장에 대한 개입은 쌍별 믿음 교환만이 제공한다.
표~\ref{tab:trace-echo}은 하나의 예측 시점을 처음부터 끝까지 보여주며,
반복된 믿음과 교정된 믿음이 생성된 추론 흔적에서 어떻게 드러나는지를 제시한다.

\begin{table*}[t]
\footnotesize
\setlength{\parindent}{0pt}
\textbf{조리 예측 사례.} 관측 구간
\obs{$3620.4$--$3628.4$}\,초, 목표 시점 \obs{$3631.0$}\,초,
가용 이력의 완료 행동 \obs{$122$}개.\par
\smallskip
\textbf{WM 후보 집합} (무작위 순서, 점수와 순위는 정책에 비공개):
check heat (정답, WM \obs{$3$}위) \textperiodcentered{} flip flatbread
\textperiodcentered{} transfer flatbread \textperiodcentered{} cook flatbread
\textperiodcentered{} press dough \textperiodcentered{} roll dough
\textperiodcentered{} knead dough (WM \obs{$1$}위) \textperiodcentered{}
sprinkle flour \textperiodcentered{} shape dough \textperiodcentered{} cut
dough.\par
\smallskip
\textbf{완료 행동 이력} (최근 \obs{$8$}개): \ldots{} $\rightarrow$ roll
dough $\rightarrow$ flip ingredient $\rightarrow$ check heat $\rightarrow$ roll
dough $\rightarrow$ check heat $\rightarrow$ flip ingredient $\rightarrow$ press
dough $\rightarrow$ cook flatbread.\par
\medskip
\setlength{\tabcolsep}{6pt}
\begin{tabular}{@{}p{0.115\textwidth}p{0.845\textwidth}@{}}
\toprule
모델 & 추론 $\rightarrow$ 믿음 $\rightarrow$ 행동 \\
\midrule
기본
  & \emph{``The current frame \scmark{shows} the flatbread being cooked on the
    hot surface, and the person's hand is positioned to flip it\ldots{}
    Flipping the flatbread is the logical next step to ensure even cooking.''}
    \newline
    \textbf{믿음} --- Flipping the flatbread to cook the other side.
    \newline
    $\rightarrow$ \textbf{flip flatbread} $\times$
    \quad{\scriptsize[``shows''를 통한 장면 묘사]} \\
GT-only
  & \emph{``The current frame \scmark{shows} the flatbread being cooked on the
    hot surface\ldots{} The most logical next step, given the context of
    cooking flatbread, is to flip it\ldots{} This aligns with
    \underline{the previous actions} of checking heat and flipping the
    flatbread.''}
    \newline
    \textbf{믿음} --- \textbf{flip flatbread}
    \quad{\scriptsize[행동을 복사한 믿음]}
    \newline
    $\rightarrow$ \textbf{flip flatbread} $\times$
    \quad{\scriptsize[``shows''를 통한 장면 묘사]} \\
후보-CE
  & \emph{``The current frame \scmark{shows} the flatbread being cooked on the
    hot surface, and \underline{the previous action} was flipping the
    ingredient. The next logical step is to continue cooking the flatbread, as
    it is still on the fire and needs to be cooked through.''}
    \newline
    \textbf{믿음} --- Cook the flatbread on the fire.
    \newline
    $\rightarrow$ \textbf{cook flatbread} $\times$
    \quad{\scriptsize[``shows''를 통한 장면 묘사]} \\
\method{}
  & \emph{``\textbf{I have just finished pressing the dough and am now
    observing the fire to assess the heat for cooking. The flames are visible
    and active, indicating the stove is ready.} I need to ensure the heat is
    suitable before proceeding with the next step in cooking.''}
    \newline
    \textbf{믿음} --- I am in the middle of assessing the heat level for
    cooking the flatbread.
    \newline
    $\rightarrow$ \textbf{check heat} \cmark
    \quad{\scriptsize[상태를 담은 믿음, 이전 행동에 대한 명시적 언급 없음]} \\
\bottomrule
\end{tabular}
\caption{\textbf{상태를 담은 믿음은 교정을 가능하게 하지만, 행동을 복사한
믿음은 그렇지 못하다.} 평가한 네 정책은 동일한 관측, 완료 행동 이력, 무작위
순서의 후보 집합, 프롬프트를 받는다. 생성된 추론 흔적은 원문 그대로 옮기며,
생략한 부분은 \ldots{}로 표시한다. 일화적 선택을 피하기 위해 공통 커버 집합에
고정된 기준을 적용한다. 네 추론 흔적이 모두 유효해야 하고, GT-only는 선택한
행동을 믿음 필드에 그대로 복사하면서 오답이어야 하며, \method{}는 정답이어야
한다. 이 기준을 만족하는 지점은 \obs{$915$}개 중 \obs{$31$}개이며, 제시한
사례는 후보-CE도 오답인 \obs{$18$}개 중 하나이다.
\underline{밑줄}은 언급률 분석과 동일한 규칙으로 이전 행동을 명시적으로 참조한
구간이고, \textbf{굵은 글씨}는 모델이 현재 과업 상태를 표현한 구간이다.
추론 흔적은 모델이 영어로 생성한 원문이므로 번역하지 않는다.}
\label{tab:trace-echo}
\end{table*}

\subsubsection{확립된 것과 확립되지 않은 것}
\label{sec:er-scope}

현재 동일 코호트의 증거는 세 가지 주장을 확립한다. 첫째, 후보 제약 학습은 월드
모델 Top-\obs{$1$} 정책과 정답 전용 감독 모두를 넘어 선택을 개선한다. 둘째,
개선은 하나의 쉬운 부분집합에만 나타나는 것이 아니라 유지와 교정 모두에서
나타난다. 셋째, 이러한 이득은 완료 행동 이력에 인과적으로 의존한다. 이 결과들은
논문의 모듈식 설명을 지지한다. 즉, 월드 모델은 시각적으로 근거가 있는 행동
경계를 정의하고, 언어 정책은 그 경계 안에서 과업과 궤적 맥락을 해석하는 법을
학습한다.

최종 학습 단계는 Retrospection-with-Replay에 근거화 오버샘플링을 더한다.
학습 집합 크기를 고정한 채 근거화 예시의 비율을 \obs{$12.8\%$}에서
\obs{$32.0\%$}로 높여 최적화 예산을 보존한다. 그 결과 \method{} 모델은
SelAcc \obs{$28.8\%$}, $G_1$ \obs{$42.5\%$}, $G_2$ \obs{$25.6\%$},
연속 행동 재현율 \obs{$48.4\%$}, 증거 유용성 \obs{$+18.3\pp$}, 믿음 반복
\obs{$0.0\%$}에 도달하며, 직전 Retrospection-with-Replay 단계에서 잃었던
선택 정확도(\obs{$26.4\%$})를 대부분 회복한다. 그럼에도
$\theta_{\mathrm{CE}}$ 대비 사전등록된 \obs{$1\pp$} 비열화 마진은 충족하지
못한다. 공식 쌍별 집합(\obs{$n{=}937$}, 비디오 클러스터 86개)에서
SelAcc 변화는 \obs{$-0.64\pp$}(\obs{$95\%$ 신뢰구간
$[-3.88,+2.91]$}), $G_2$ 변화는 \obs{$-2.53\pp$}
(\obs{$95\%$ 신뢰구간 $[-6.98,+2.29]$})이다. 두 구간 모두 0을 포함하며
통상적인 \obs{$4\pp$} 마진이라면 충족하겠지만, 그러한 해석은 사후적이다.
따라서 Retrospection이
Prospection을 성능 저하 없이 보존한다고 주장하지 않고 최종 모델을 보고한다.

{\color{pendingred}
\paragraph{최종 단계 개입의 한계.}
근거화 오버샘플링 이후 믿음 교환 평가를 수행하지 않았으므로, 최종 \method{}
모델에 대해 축 5는 확립되지 않았다. 직전 Retrospection-with-Replay 모델의
믿음 교환 민감도는 $0.098$로 사전등록된 $0.20$ 임계값보다 낮았던 반면,
추론 교환 민감도는 $0.483$으로 평가한 조건 가운데 가장 높았다.
즉 개입 질량이 사라진 것이 아니라 추론 채널로 이동했다.
최종 믿음--행동 반복률은 $0.0\%$이지만, 반복률만으로는 개입 검정이 되지
않는다. 따라서 긍정적인 믿음 유용성 진단은 탐색적 관찰로만 유지한다.
\par}

{\color{pendingred}
\paragraph{제외한 대리 지표.}
1인칭 표현은 의도적으로 역량 프로파일에서 제외한다. GT-only 대조군에서는 후보가
제시된 경우 $30.7\%$였던 비율이 후보 없는 생성에서는 $91.0\%$로
변하며, GT-only 추론 흔적 내에서도 정확도를 양으로 예측하지 않는다
(1인칭 시작 $16.1\%$ 대 주석 스타일 시작 $29.7\%$).
따라서 이는 체화 추론의 증거가 아니라 프롬프트와 템플릿에 민감한 표면 형식이다.
대조 접속 표현도 마찬가지로 제외한다. 기본 모델 추론 흔적의 $0.3\%$,
최종 모델 추론 흔적의 $1.9\%$에서만 나타나며, 조건 내 정확도 연관성은
재생 단계와 근거화 단계 사이에서 불안정하다(각각 $+19.2\pp$와
$-7.6\pp$, 최종 단계 추론 흔적은 $18$건뿐). 이는 소표본 잡음과 일관된다.
어느 대리 지표도 논문의 주장을 뒷받침하는 데 사용하지 않는다.
\par}

\end{document}
